

\documentclass[tikz,svgnames,border={10 10},convert={density=300, outext=.png}]{standalone}

% 若需要中文，请用 xelatex 编译并取消下面两行注释：
\usepackage{xeCJK}
% \setCJKmainfont{Noto Sans CJK SC}

\usetikzlibrary{positioning,arrows,fit}

\begin{document}
\begin{tikzpicture}[
    box/.style={rectangle,draw,fill=DarkGray!20,rounded corners,
                node distance=1cm,text width=18em,text centered,
                minimum height=2.4em,thick},
    arrow/.style={draw,-latex',thick},
  ]

  % === 主体流程节点（与前述模块一致风格，不加外框） ===
  \node[box,fill=PeachPuff!60,text width=24em] (input)
    {接收配置与命令行输入：YAML/JSON 配置、CLI 覆盖、参数扫描定义};

  \node[box,below=0.6 of input] (parse)
    {配置解析与类型验证（parse / override → 统一配置对象）};

  \node[box,below=0.5 of parse] (sweep)
    {参数扫描与调度（generate\_experiments / submit\_jobs）};

  \node[box,below=0.5 of sweep] (setup)
    {初始化仿真环境（setup\_experiment / 依赖锁 / 随机种子）};

  \node[box,below=0.5 of setup] (run)
    {运行与监控（run；日志/状态机/失败恢复）};

  \node[box,below=0.5 of run] (crit)
    {还有待运行的实验？};

  % 底部输出
  \node[box,below=1.8 of crit,fill=LightSteelBlue!60,text width=26em] (output)
    {输出：实验运行上下文、日志目录结构、结果索引与归档包；配置快照与锁文件};

  % === 箭头（直连 + Yes/No 回路） ===
  \path[arrow] (input) -- (parse);
  \path[arrow] (parse) -- (sweep);
  \path[arrow] (sweep) -- (setup);
  \path[arrow] (setup) -- (run);
  \path[arrow] (run) -- (crit);

  % Yes 分支：直达输出
  \path[arrow] (crit) -- (output)
    node[midway,left=0.1,draw,outer sep=0pt,fill=white] {No};

  % No 分支：先竖直向下，再水平右移（Yes 在水平段下方），最后折返到“初始化仿真环境”节点
  % 语义：若“还有待运行的实验？”为 Yes，则回到 setup 继续下一轮
  \path[draw,thick] (crit.south) ++(0,-1em) coordinate(novert);
  \draw[thick] (novert) -- ++(4,0) coordinate(nojump)
        node[midway,below=2pt,draw,outer sep=0pt,fill=white] {Yes};
  \draw[arrow] (nojump) |- (setup.east);

\end{tikzpicture}

\end{document}